\par\medskip
\begin{center}d) \textls[100]{Anwendungen. Zusammenfassung.}\end{center}
\noindent\hspace*{\parindent}Aus der Tatsache, dass \(k'=q(t)^2/p(t)^2\) in dem oben charakterisierten Bereiche
der rechten \(t\)-Halbebene jeden komplexen Wert annimmt, folgt nunmehr
ohne weiteres, dass der Ansatz (11), bezw. (11a) für jeden komplexen
Wert von \(k'\), also für jedes komplexe Wertepaar \(a,b\) möglich ist, was eine
in sich abgeschlossene Theorie des agM. in allgemeinster Fassung liefert. Eine
dieser »transzendenten« Theorie entsprechende algebraische, d. h. von dem Algorithmus
ausgehende Theorie hat \textsc{Gauss} nicht aufgezeichnet\footnote{Eine solche hat L. v. \textsc{Dávid}, \textsc{Crelles} Journal 159, 1928, S. 154 skizziert, vergl. auch \textsc{Geppert}, Math. Annalen 99, 1928, S. 162 und Jahresbericht der D.M.-V. 38, 1929, S. 73.}; er hat aber
Wert darauf gelegt, für die Berechnung der zu einem gegebenen Wertepaar
\(a,b\) oder einem gegebenen Werte \(k'\) gehörigen Grösse \(x=e^{-\pi t}\) brauchbare
Formeln zu geben. Es finden sich solche an verschiedenen Stellen des Nachlasses,
zuerst wohl im Art. 2 der Aufzeichnung 2) (1801, Werke X, 1, S. 210),
dann im Art. 2 der Aufzeichnung 3) (ebenda S. 218), im Art. 11 derselben
Aufzeichnung (ebenda S. 227), ferner in der aus 1809 stammenden Notiz VIII\par
\clearpage\SourceLeaf{192}\ArithmeticHead{107}{R}{Vd. ANWENDUNGEN. ZUSAMMENFASSUNG.}
\HistoricalFoot{}{14\textsuperscript{*}}
\noindent (ebenda S. 213) und in den \textit{Hundert Theoremen}, Art. 6 (Werke III, S. 467,
468). Er verfährt dabei so, dass er die Transformation (9) oder, was dasselbe
heisst, den Übergang von \(t\) zu \(2t\) unendlich oft anwendet und dadurch für \(x\)
die Darstellung erhält:
\[
\begin{aligned}
x&=\lim_{n\to\infty}\left(\tfrac12r(2^nt)\right)^{\frac1{2^{n-1}}}
=\lim_{n\to\infty}\left(\frac{c_n}{4a_n}\right)^{\frac1{2^{n-1}}}\\
&=\lim_{n\to\infty}(\tfrac14k_n)^{\frac1{2^{n-1}}},
\end{aligned}
\]
\noindent wo (vergl. die Gl. (9a) oben S. 92)
\[
k_n=\frac{1-k'_{n-1}}{1+k'_{n-1}},\quad k'_{n-1}=\sqrt{1-k_{n-1}^2},
\]
\noindent ist, oder in Gestalt eines unendlichen Produkts
\[
x=\frac{a^2-b^2}{16a_1a_2}\left(\frac{a_2}{a_3}\right)^{\frac12}\left(\frac{a_3}{a_4}\right)^{\frac14}\ldots.
\]
\noindent »Da die Grössen \(a_1,a_2,a_3,\ldots\) sich äusserst schnell der Gleichheit nähern«,
sagt \textsc{Gauss} (Werke III, S. 469), »erhält man \(x\) mit grösster Bequemlichkeit«.
Zieht man die Formel (8) Abschn. IV (oben S. 67),
\[
\lim M(1,p)\log\frac4p=\frac12\pi
\]
\noindent oder eine der mit ihr äquivalenten\footnote{Etwa die Gl. (10) der Aufzeichnung 2), Werke X, 1, S. 211, oder den »Schönen Lehrsatz« der Aufzeichnung 3) Art. 4, ebenda, S. 220, oder das »Theorema elegantissimum« Werke VIII, S. 98.} heran, so kann man von dieser Darstellung
von \(x\) wieder zu der Formel (18) oder
\[
x=e^{-\pi\frac{M(a,b)}{M(a,c)}}
\]
\noindent gelangen\footnote{Vergl. Art. 2 der Aufzeichnung 2), Werke X, 1, S. 210, Notiz VIII, ebenda S. 213 und Art. 11 der Notiz 3), ebenda S. 227; siehe auch \textsc{Geppert} a. a. O. \textsuperscript{83)} § 25, S. 104.}.

Wenn wir uns nun der Anwendung der geschilderten Theorie auf die
eigentliche Theorie der elliptischen Funktionen zuwenden, so können wir zunächst
die Frage der Umkehrung eines Integrals erster Gattung (1) (oben
S. 83) mit einem beliebig vorgeschriebenen \(\mu^2\) jetzt mit wenigen Worten erledigen.
Aus der Gleichung (vergl. (3), oben S. 86): \(\mu^2=(1-k'^2)/k'^2\) bestimmen
wir den zu \(\mu^2\) gehörigen Wert von \(k'^2\); dann gibt es im Fundamentalbereich \(F_2\)
ein und nur ein \(t\), für das \(q(t)^4/p(t)^4\) diesen Wert \(k'^2\) annimmt, und wir kennen\par
\clearpage\SourceLeaf{193}\ArithmeticHead{108}{L}{L. SCHLESINGER, ÜBER GAUSS’ FUNKTIONENTHEORETISCHE ARBEITEN.}
\noindent einerseits rechnerisch gut brauchbare Methoden, um diesen Wert \(t\) oder das
entsprechende \(x\) zu bestimmen, und können andererseits auch alle \(t\)-Werte
charakterisieren, die zu diesem \(k'^2\) führen. Die mit einem dieser \(t\) gebildete
Funktion \(\operatorname{Sl}\varphi\) löst dann das in Rede stehende Umkehrproblem.

Von weiteren Anwendungen berichten zwei Tagebuchnotizen, die in unmittelbarem
Anschluss an die bisher besprochenen niedergeschrieben sind und
wie folgt lauten: 110. vom 5. Juni 1800, »Unsere Theorie haben wir auch
unmittelbar auf die elliptischen Transzendenten angewandt« und 111. vom
10. Juni 1800, »Die Rektifikation der Ellipse wurde auf drei verschiedene
Weisen erledigt«. Dazu ist zu bemerken, dass »elliptische Transzendente« bei
\textsc{Gauss} die Bedeutung von dem hat, was wir elliptische Integrale erster Gattung
nennen; so spricht er in der oben (S. 71) wiedergegebenen Aufzeichnung in
bezug auf das Integral \(\displaystyle\int\frac{d\varphi}{\sqrt{1-e^2\sin^2\varphi}}\) auch von der \textls[100]{elliptischen Funktion}.
Das vollständige bezw. allgemeine Integral \textls[100]{zweiter} Gattung bezeichnet \textsc{Gauss}
stets als den Ellipsenquadranten bezw. als Rektifikation der Ellipse, so dass
also die Notiz 111. auf das Integral zweiter Gattung mit variabler oberer
Grenze zu deuten ist. In diesem Sinne finden wir die Ausführung der am
5. und 10. Juni notierten Untersuchungen in den Artikeln 10—14 der Notiz 3)
Werke X, 1, S. 227—231, von denen die Artikel 13, 14 direkt »Zur Rektifikation
der Ellipse« überschrieben sind.

\textsc{Gauss} bedient sich hier der beiden Transformationen, deren Darstellung
durch die Thetafunktionen in der Theorie des allgemeinen Sinus lemniscaticus
der Scheda Ac (Artikel 1—3, Werke X, 1, S. 194—196) enthalten
ist, und die wir schon oben bei der Wiedergabe der betreffenden Formeln
(S. 78) als \textsc{Landensche} bezw. \textsc{Gausssche} Transformation bezeichnet haben. —
Für die letztere wird sogar im Art. 8 der Scheda Ac-Aufzeichnung (Werke
X, 1, S. 200) die Form angegeben, in der sie z. B. im Art. 16 der Abhandlung
\textit{Determinatio attractionis} von 1818 (Werke III, S. 352) auftritt, so
dass es sehr wahrscheinlich ist, dass \textsc{Gauss} diese Formel aus den Gleichungen
für die Transformation der Thetafunktionen, wie wir sie im Art. 3 finden,
abgeleitet hat. Wenn man nämlich unsere Gleichungen (27a), (27b) Abschnitt
IV (oben S. 78) durcheinander dividiert und\footnote{Vergl. die Fussnote \textsuperscript{115)} und die daselbst zitierte Notiz Werke III, S. 473, Art. 5.}
\clearpage\SourceLeaf{194}\ArithmeticHead{109}{R}{Vd. ANWENDUNGEN. ZUSAMMENFASSUNG.}
\[
\begin{gathered}
\frac{p(x^{\frac12})}{r(x^{\frac12})}\frac{\vartheta_{10}(\psi|x^{\frac12})}{\vartheta_{00}(\psi|x^{\frac12})}
=\operatorname{Cl}\varphi=\sin T,\\
\varphi=\pi\psi q(x^{\frac12})^2,\quad\operatorname{tg}v=\frac{r(x^{\frac12})^2}{q(x^{\frac12})^2},\\[5pt]
\frac{p(x)}{r(x)}\frac{\vartheta_{10}(\psi|x)}{\vartheta_{00}(\psi|x)}
=\operatorname{Cl}\varphi_1=\sin T_1,\\
\varphi_1=\pi\psi q(x)^2,\quad\operatorname{tg}v_1=\frac{r(x)^2}{q(x)^2}
\end{gathered}
\]
\noindent setzt, so dass (vergl. (9), oben S. 92) \(\sin v_1=\operatorname{tg}^2\frac12v\) wird\footnote{\textsc{Gauss} bezeichnet Werke X, 1, S. 200 die dem \(x\) entsprechenden Grössen mit \(\varphi,v\), die dem \(x^{\frac12}\) entsprechenden mit \({}'\varphi,{}'v\), womit er offenbar andeuten will, dass es sich beim Übergang von \(x\) zu \(x^{\frac12}\) um eine Fortsetzung nach der negativen Seite hin handelt. Auch schreibt er die Beziehung nicht, wie wir es im Text tun, für Cl, sondern für Sl, was natürlich ganz irrelevant ist.}, so erhält man
\[
\sin T=\frac{\sin T_1}{\frac{1+\cos v}2+\frac{1-\cos v}2\sin^2T_1},
\]
\noindent oder, indem man \(\cos v=b/a\) setzt, die in der \textit{Determinatio attractionis} (Werke
III, S. 352) auftretende Form
\begingroup\makeatletter\tagsleft@true\let\veqno\leqno\makeatother
\begin{equation*}
\sin T=\frac{2a\sin T_1}{a+b+(a-b)\sin^2T_1},\tag{26}
\end{equation*}
\endgroup
\noindent für die dann (vergl. Werke X, 1, S. 280, aus Handbuch 16, Bb) die Gleichung
\[
\begin{gathered}
\frac{dT}{\sqrt{a^2\sin^2T+b^2\cos^2T}}=
\frac{dT_1}{\sqrt{a_1^2\sin^2T_1+b_1^2\cos^2T_1}},\\
a_1=\tfrac12(a+b),\quad b_1=(ab)^{\frac12},
\end{gathered}
\]
\noindent gilt. Dies gäbe ins Unendliche fortgesetzt schon die Methode zur Berechnung
des Integrals
\begingroup\makeatletter\tagsleft@true\let\veqno\leqno\makeatother
\begin{equation*}
\int\frac{dT}{\sqrt{a^2\sin^2T+b^2\cos^2T}}=\frac{\Theta}{M(a,b)},\tag{27}
\end{equation*}
\endgroup
\noindent die wir im Art. 18 der \textit{Determ. attr.} finden, die aber numerisch nicht zweckmässig
ist. Deshalb nimmt \textsc{Gauss} die \textsc{Landensche} Transformation (Gln. (24),
(25), oben S. 78) zu Hilfe und setzt (Artikel 10, 12, Werke X, 1, S. 227—229)
\[
\begin{gathered}
\operatorname{tg}2^{n-1}T_n=\frac{\vartheta_{11}(2^n\psi|x^{2^n})}{\vartheta_{00}(2^n\psi|x^{2^n})},\\
\operatorname{tg}2^nU_n=\frac{p(x^{2^{n+1}})}{q(x^{2^{n+1}})}
\frac{\vartheta_{11}(2^{n+1}\psi|x^{2^{n+1}})}{\vartheta_{00}(2^{n+1}\psi|x^{2^{n+1}})},\\
\psi=\frac{\Theta}{2\pi},\quad x=e^{-\pi\frac{M(a,b)}{M(a,c)}},\qquad(n=0,1,2\ldots)
\end{gathered}
\]
\noindent wo dann \(U_n=\frac12(U_{n-1}+T_{n-1})\) ist. Die rekursive Berechnung der \(T_n,U_n\) führt
auf trigonometrische Algorithmen, die für logarithmische Rechnung sehr praktikabel
sind, und die für das Integral erster Gattung (27) direkt \(\Theta=\lim\limits_n T_n=\)\par
\clearpage\SourceLeaf{195}\ArithmeticHead{110}{L}{L. SCHLESINGER, ÜBER GAUSS’ FUNKTIONENTHEORETISCHE ARBEITEN.}
\noindent \(\lim\limits_n U_n\) ergeben, während sie für das Integral zweiter Gattung
\[
\int\sqrt{(a^2\sin^2T+b^2\cos^2T)}\,dT
\]
\noindent eine unendliche Reihe liefern. — Für dieses letztere Integral wird im Art. 13
eine zweite Methode entwickelt, die auf der unendlich oftmaligen Wiederholung
der \textsc{Gaussschen} Transformation (27) Abschn. IV (oben S. 78) beruht,
die durch Einführung passender Parameter zu einem für die Rechnung bequemen
Algorithmus gestaltet wird\footnote{Vergl. die Darstellung der Berechnungen für die Integrale erster und zweiter Gattung bei \textsc{Geppert} a. a. O. \textsuperscript{83)}, §§ 40, 41, S. 167—177.}.

Von der in der Tagebuchnotiz erwähnten dritten Methode für die Rektifikation
hat sich im Nachlass nichts vorgefunden; vielleicht hat \textsc{Gauss} auch
die Darstellung des Art. 13 als \textls[100]{zwei} Methoden gezählt, weil sie zwei äusserlich
verschiedene Reihenentwicklungen neben einander gibt. — In bezug auf
die Darstellung des Integrals erster Gattung wird die weiter unten folgende
Besprechung des Nachlassstücks VIII (Werke X, 1, S. 213ff.), in bezug auf
die des Integrals zweiter Gattung die \textit{Determinatio attractionis}, Art. 16 (Werke
III, S. 353) und die \textit{Anzeige} dieser Abhandlung (Werke III, S. 359, 360) zu
vergleichen sein. Was die Abfassungszeit der Art. 10—14 der Notiz 3) anlangt,
die drei Zetteln des Nachlasses entnommen sind, so dürften sie (vergl.
die Bemerkungen Werke X, 1, S. 283) aus dem zweiten Jahrzehnt des 19.
Jahrhunderts stammen.

\textsc{Gauss} scheint dann längere Zeit hindurch die jetzt zu einem gewissen
Abschluss gebrachte Theorie der elliptischen Funktionen nicht weiter verfolgt
zu haben. Wir werden sehen, dass er seinen hier geschilderten Untersuchungen
auch später keine wesentlich neuen Gedanken hinzugefügt hat, dass er namentlich
immer an der Theorie des agM. und der Thetafunktionen haften blieb
und eine direkte funktionentheoretische Untersuchung weder für das Integral
erster Gattung noch für seine Umkehrungsfunktion angestellt hat. In bezug
auf diese letztere Untersuchung hat er also den Arbeiten von \textsc{Jacobi}, \textsc{Abel},
\textsc{Puiseux} und \textsc{Riemann}, die dann weiter zur Theorie der Integrale allgemeiner
algebraischer Funktionen und ihrer Umkehrung führten, nicht vorgegriffen.
Dagegen hat \textsc{Gauss} die Theorie der elliptischen Funktionen sogleich auf eine
breitere Basis gestellt, indem er seinen allgemeinen Sinus lemniscaticus und\par
\clearpage\SourceLeaf{196}\ArithmeticHead{111}{R}{Vd. ANWENDUNGEN. ZUSAMMENFASSUNG.}
\noindent die vier Thetafunktionen \(P,Q,R,S\) von vornherein als Funktionen der \textls[100]{beiden}
unabhängigen Variablen \(x=e^{-\pi t},\ y=e^{2\psi\pi i}\) oder \(t\) und \(\psi\) betrachtet.
Dass ihn dabei die Abhängigkeit von \(x\), also von \(t\), ganz besonders interessiert
hat, kommt ja schon in der Reihenfolge zur Geltung, in der er die beiden
Variabeln schreibt\footnote{Dass eine solche Äusserlichkeit bei \textsc{Gauss} auch stets die Folge besonderer Überlegung war, wird z. B. durch eine Mitteilung bestätigt, die \textsc{Encke} in seiner Erwiderung auf die Antrittsrede \textsc{Kroneckers} in der Berliner Akademie (Monatsberichte 1861, \textsc{Kroneckers} Werke V, 1930, S. 391) macht. Er sagt: »Mein hochgeehrter Lehrer .\ .\ . \textsc{Gauss} .\ .\ . pflegte in vertraulichem Gespräch häufig zu äussern, die Mathematik sei weit mehr eine Wissenschaft für das Auge, als eine für das Ohr .\ .\ .\ . Eben deshalb war \textsc{Gauss} in der Wahl der Bezeichnungen, selbst bei der Auswahl der einzelnen Buchstaben, ungemein peinlich .\ .\ .«}; in der Tat ist er gerade in bezug auf die Untersuchung
dieser Abhängigkeit nicht nur, wie schon M. \textsc{Noether} bemerkt hat\footnote{M. \textsc{Noether}, Mathem. Annalen 55, 1902, S. 370.}, \textsc{Abel}
und \textsc{Jacobi}, sondern auch allen späteren Bearbeitern der Theorie der elliptischen
Funktionen von \textsc{Cauchy-Liouville} bis \textsc{Weierstrass} dauernd überlegen
geblieben. Die ihm aus früher Jugendzeit geläufige Methode des Operierens
mit Reihen, besonders Potenzreihen, hat er verfeinert und in ihrer Kraft dadurch
ausserordentlich gesteigert, dass er sie mit dem Studium der durch
solche Reihen vermittelten Abbildung, also mit den Hilfsmitteln der geometrischen
Funktionentheorie verknüpfte, Hilfsmittel, deren er sich ja schon in
der Inauguraldissertation mit Erfolg zu bedienen wusste, und deren Anwendungen
uns dann in jener Zettelaufzeichnung über das Integral erster Gattung
(oben S. 71) und in den Notizen der Scheda An (oben S. 102) begegnen. So
hat er es vermocht, um die Wende vom 18. zum 19. Jahrhundert Wege zu
beschreiten, denen die Mathematiker auch zur Zeit, als \textsc{Schering} im Bande III
der Werke Nachlassstücke veröffentlichte (1868), noch nicht zu folgen imstande
waren. Immerhin hat aber jene Veröffentlichung \textsc{Gauss} auch die Priorität
der Publikation gesichert für die oben von uns skizzierte Theorie der Modulfunktion,
namentlich bleibt z. B. die Tatsache bestehen, dass im Jahre 1868
auf S. 477 von Werke III der Fundamentalbereich der Modulfunktion zwar
undeutlich gezeichnet, aber arithmetisch durch die Bedingung, dass der imaginäre
Teil von \(t\) und \(1/t\) zwischen \(+i\) und \(-i\) gelegen sei, einwandfrei charakterisiert
und in seiner vollen Bedeutung erfasst vorlag. Dass \textsc{Dedekind}
beim Aufbau seiner Theorie der Modulfunktion an die Bemerkung Werke III,
S. 368 angeknüpft hat, wurde schon oben bemerkt.
